\documentclass[12pt]{article}
\usepackage[utf8]{inputenc}
\usepackage[T1]{fontenc}
\usepackage[margin=1in]{geometry}
\usepackage{setspace}
\usepackage{xcolor}
\usepackage[colorlinks=true, urlcolor=blue!60!black, linkcolor=blue!60!black]{hyperref}
\setlength{\parindent}{2em}
\setlength{\parskip}{1.25em}
\renewcommand{\baselinestretch}{1.0}
\begin{document}
\begin{center}
{\LARGE \textit{United States v. Carolene Products Co.}, 304 U.S. 144 (1938)\par}
\vspace{0.8em}
{\large Case Brief\par}
\end{center}
\vspace{1em}

\section{Memory Jogger}

In upholding Congress's ban on interstate shipments of ``filled milk,'' the Supreme Court cemented deferential rational basis review for ordinary economic regulations while planting in Footnote Four the doctrinal blueprint for heightened judicial scrutiny of laws infringing specific constitutional rights, distorting the political process, or targeting ``discrete and insular minorities.''

\section{Detailed Case Facts}

On March 4, 1923, Congress enacted the Filled Milk Act, ch.~262, 42 Stat.~1486 (codified at 21 U.S.C.~\S\S~61--63), following extensive committee hearings before the House Committee on Agriculture (H.R.~Rep.~No.~365, 67th Cong., 1st Sess.) and the Senate Committee on Agriculture and Forestry (S.~Rep.~No.~987, 67th Cong., 4th Sess.) (\href{https://supreme.justia.com/cases/federal/us/304/144/}{304 U.S. at 145--49}; \href{https://www.courtlistener.com/opinion/103017/united-states-v-carolene-products-co/}{CourtListener}). Section~61(c) defined ``filled milk'' as any milk, cream, or skimmed milk---whether or not condensed, evaporated, concentrated, powdered, dried, or desiccated---blended or compounded with any fat or oil other than milk fat so that the resulting product was ``in imitation or semblance of milk, cream, or skimmed milk'' (\href{https://supreme.justia.com/cases/federal/us/304/144/}{304 U.S. at 145--46 n.1}). Section~62 declared filled milk to be ``an adulterated article of food, injurious to the public health, and its sale constitutes a fraud upon the public,'' and made it unlawful to ship or deliver it for shipment in interstate or foreign commerce (\href{https://supreme.justia.com/cases/federal/us/304/144/}{304 U.S. at 145--46 n.1}). Section~63 imposed criminal penalties of a fine up to \$1,000, imprisonment up to one year, or both.

Congressional committee reports concluded that extracting natural butterfat---a vital source of vitamins, particularly for children---and substituting cheaper vegetable oils such as coconut oil created a lower-cost product that looked and tasted like pure condensed milk, invited deceptive marketing despite labeling laws under the Pure Food and Drugs Act, and risked malnutrition (\href{https://supreme.justia.com/cases/federal/us/304/144/}{304 U.S. at 149 n.2}). By 1938, roughly thirty-five states had enacted statutes prohibiting filled milk, and three more heavily regulated it (\href{https://supreme.justia.com/cases/federal/us/304/144/}{304 U.S. at 150 n.3}).

The appellee, Carolene Products Company, manufactured and shipped in interstate commerce a filled-milk product called ``Milnut,'' a compound of condensed skimmed milk and coconut oil made in imitation or semblance of condensed milk or cream (\href{https://supreme.justia.com/cases/federal/us/304/144/}{304 U.S. at 146}). Federal prosecutors indicted the company for violating the Filled Milk Act, alleging in the statutory language that Milnut was an adulterated food injurious to public health and did not fall within any statutory exception. Carolene Products challenged the Act on its face, contending that it exceeded Congress's power under the Commerce Clause (U.S. Const.~art.~I, \S~8, cl.~3), invaded state authority reserved by the Tenth Amendment, violated the Fifth Amendment's Due Process Clause by making the legislative declaration of harm conclusive against the company, and arbitrarily singled out filled milk while leaving oleomargarine (a butter substitute using vegetable oil) unregulated (\href{https://supreme.justia.com/cases/federal/us/304/144/}{304 U.S. at 146--47, 151}).

\section{Procedural History}

\begin{itemize}
    \item \textbf{Prior District Court Ruling (1934):} In an earlier prosecution involving the same company, the U.S.~District Court for the Southern District of Illinois held the Filled Milk Act unconstitutional in \href{https://law.justia.com/cases/federal/district-courts/FSupp/7/500/1943047/}{\textit{United States v. Carolene Products Co.}, 7 F.~Supp.~500 (S.D.~Ill.~1934)}.
    \item \textbf{District Court Ruling in This Case:} Following a new indictment in the U.S.~District Court for the Southern District of Illinois for interstate shipment of Milnut, Carolene Products filed a demurrer challenging the facial constitutionality of the Filled Milk Act. Relying on its 1934 precedent (\textit{7 F.~Supp.~500}), the district court sustained the demurrer and dismissed the indictment (\href{https://supreme.justia.com/cases/federal/us/304/144/}{304 U.S. at 146}).
    \item \textbf{Parallel Circuit Authority:} Meanwhile, in separate civil litigation, the U.S.~Court of Appeals for the Seventh Circuit upheld the Filled Milk Act as a valid exercise of the commerce power in \textit{Carolene Products Co. v. Evaporated Milk Ass'n}, 93 F.2d 202 (\href{https://supreme.justia.com/cases/federal/us/304/144/}{304 U.S. at 146}).
    \item \textbf{Direct Appeal to the U.S.~Supreme Court:} The United States appealed directly from the district court to the Supreme Court under the Criminal Appeals Act of March 2, 1907, 34 Stat.~1246, 18 U.S.C.~\S~682. The case was argued on April 6, 1938, and decided on April 25, 1938. The Supreme Court reversed, holding that the demurrer should have been overruled (\href{https://supreme.justia.com/cases/federal/us/304/144/}{304 U.S. at 145, 154}). (In subsequent litigation following a later conviction, the Court again upheld the Act against an as-applied Fifth Amendment challenge in \href{https://supreme.justia.com/cases/federal/us/323/18/}{\textit{Carolene Products Co. v. United States}, 323 U.S. 18 (1944)}.)
\end{itemize}

\section{Judicial Votes}

The Court voted \textbf{6--1 to reverse} the judgment sustaining the demurrer, with seven Justices participating (\href{https://supreme.justia.com/cases/federal/us/304/144/}{304 U.S. at 154--55}; \href{https://constitutioncenter.org/the-constitution/supreme-court-case-library/united-states-v-carolene-products-co}{National Constitution Center}):

\begin{itemize}
    \item \textbf{Majority Opinion (Parts First, Second, and Third, including Footnote Four --- 4 Justices):} Justice Harlan Fiske Stone (author), joined in full by Chief Justice Charles Evans Hughes, Justice Louis D. Brandeis, and Justice Owen J. Roberts.
    \item \textbf{Concurring in Part and in the Result (1 Justice):} Justice Hugo L. Black concurred in the result and joined all of Justice Stone's opinion \textit{except} the part marked ``Third''---meaning Justice Black joined Parts First and Second (giving them a 5-Justice majority) but withheld his vote from Part Third, which contained the rational-basis formulation, the discussion of as-applied factual challenges, and Footnote Four (\href{https://supreme.justia.com/cases/federal/us/304/144/}{304 U.S. at 155}).
    \item \textbf{Concurring in the Result Only (1 Justice):} Justice Pierce Butler filed a separate opinion concurring solely in the result, arguing that while the indictment was prima facie sufficient to survive a demurrer, the defendant must be allowed to prove at trial that Milnut was wholesome and non-deceptive (\href{https://supreme.justia.com/cases/federal/us/304/144/}{304 U.S. at 155}).
    \item \textbf{Dissenting (1 Justice):} Justice James Clark McReynolds dissented without opinion, noting only that he ``thinks that the judgment should be affirmed'' (\href{https://supreme.justia.com/cases/federal/us/304/144/}{304 U.S. at 155}).
    \item \textbf{Not Participating (2 Justices):} Justice Benjamin N. Cardozo and Justice Stanley F. Reed took no part in the consideration or decision of the case (\href{https://supreme.justia.com/cases/federal/us/304/144/}{304 U.S. at 155}).
\end{itemize}

\section{Holding}

\textbf{Judgment:} Reversed. The Filled Milk Act is not unconstitutional on its face, and the district court should have overruled the demurrer to the indictment (\href{https://supreme.justia.com/cases/federal/us/304/144/}{304 U.S. at 154}).

\textbf{Rule Announced:} Congress may prohibit the interstate shipment of articles it reasonably conceives to be injurious to public health, morals, or welfare without violating the Commerce Clause, Tenth Amendment, or Fifth Amendment Due Process Clause. Regulatory legislation affecting ordinary commercial transactions carries a strong presumption of constitutionality and will not be invalidated unless, in light of facts made known or generally assumed, it precludes the assumption that it rests upon some rational basis within the knowledge and experience of the legislators (\href{https://supreme.justia.com/cases/federal/us/304/144/}{304 U.S. at 152}). Moreover, the Fifth Amendment contains no Equal Protection Clause, and legislatures are not required to address all related evils at once (\href{https://supreme.justia.com/cases/federal/us/304/144/}{304 U.S. at 151}).

\textbf{Holding vs.~Dicta:} The narrow, binding \textbf{holding} is that the Filled Milk Act is facially valid under the Commerce Clause and Fifth Amendment Due Process Clause, so the district court erred in dismissing the indictment on demurrer. By contrast, \textbf{Footnote Four}---suggesting narrower deference or ``more exacting judicial scrutiny'' for laws that violate specific constitutional prohibitions, restrict political processes, or reflect prejudice against ``discrete and insular minorities''---is \textbf{dicta} (\href{https://supreme.justia.com/cases/federal/us/304/144/}{304 U.S. at 152--53 n.4}). Not only was Footnote Four unnecessary to resolve a facial challenge to a milk-substitute statute, but because Justice Black declined to join Part Third, Part Third and Footnote Four commanded only a four-Justice plurality (Stone, Hughes, Brandeis, and Roberts) on a seven-Justice Court (\href{https://chicagounbound.uchicago.edu/cgi/viewcontent.cgi?article=2988&context=journal_articles}{Strauss, \textit{Is Carolene Products Obsolete?}}). Similarly, Part Third's remarks preserving as-applied challenges when supporting facts cease to exist were dicta in a facial demurrer posture. Despite its technical status as plurality dicta, Footnote Four became the foundational framework of modern tiered constitutional scrutiny.

\section{Analysis of Opinions}

\subsection{Opinion of the Court (Justice Stone)}

Justice Stone structured his opinion in three numbered parts (\href{https://supreme.justia.com/cases/federal/us/304/144/}{304 U.S. at 147--54}).

In \textbf{Part First}, Stone confirmed that Congress's power under the Commerce Clause is ``complete in itself'' (\textit{Gibbons v. Ogden}) and includes the authority to exclude from interstate commerce articles that Congress reasonably conceives to be injurious to public health, morals, or welfare. Because this authority is plenary within its sphere, exercising it in ways that resemble a state police power does not invade powers reserved to the states under the Tenth Amendment (\href{https://supreme.justia.com/cases/federal/us/304/144/}{304 U.S. at 147--48}).

In \textbf{Part Second}, Stone rejected the Fifth Amendment due process and equality arguments. Relying on \textit{Hebe Co. v. Shaw}, 248 U.S. 297 (1919) (which upheld an Ohio ban on filled milk under the Fourteenth Amendment), Stone emphasized both the presumption of constitutionality and the substantial affirmative record compiled in congressional hearings showing nutritional deficits and consumer confusion. Whether labeling under the Pure Food and Drugs Act was sufficient or outright prohibition was preferable was ``a matter for the legislative judgment, and not that of courts'' (\href{https://supreme.justia.com/cases/federal/us/304/144/}{304 U.S. at 151}). Stone also dismissed the objection that Congress did not ban oleomargarine, noting that the Fifth Amendment has no Equal Protection Clause and that even under the Fourteenth Amendment, ``[a] legislature may hit at an abuse which it has found, even though it has failed to strike at another'' (\href{https://supreme.justia.com/cases/federal/us/304/144/}{304 U.S. at 151}).

In \textbf{Part Third} (joined by Hughes, Brandeis, and Roberts, but not Black), Stone addressed Carolene Products' claim that Section~62's statutory declaration of health injury and fraud unconstitutionally foreclosed judicial review. Stone treated the statutory declaration merely as a statement of legislative findings aiding informed review, no more prejudicial than committee reports or surplusage, because facts supporting ordinary commercial legislation are presumed to exist unless precluded by known or generally assumed facts (\href{https://supreme.justia.com/cases/federal/us/304/144/}{304 U.S. at 152--53}). While acknowledging that factual inquiries might be appropriate in as-applied challenges or where underlying facts change, Stone stressed that where the legislative judgment is ``at least debatable,'' neither a court nor a jury may substitute its view for Congress's (\href{https://supreme.justia.com/cases/federal/us/304/144/}{304 U.S. at 154}).

Appended to the presumption-of-constitutionality sentence in Part Third was \textbf{Footnote Four}, which articulated three categories where the presumption might have ``narrower scope'' (\href{https://supreme.justia.com/cases/federal/us/304/144/}{304 U.S. at 152--53 n.4}):
\begin{enumerate}
    \item \textit{Specific Constitutional Prohibitions:} Legislation that appears on its face to fall within a specific prohibition of the Constitution, such as the first ten amendments (including as incorporated through the Fourteenth Amendment).
    \item \textit{Political Process Restrictions:} Legislation restricting ``those political processes which can ordinarily be expected to bring about repeal of undesirable legislation,'' such as restrictions on voting, dissemination of information, political organizations, or peaceable assembly.
    \item \textit{Prejudice Against Discrete and Insular Minorities:} Statutes directed at particular religious, national, or racial minorities, where ``prejudice against discrete and insular minorities may be a special condition'' curtailing ordinary political protections and warranting ``more searching judicial inquiry.''
\end{enumerate}

\subsection{Partial Concurrence (Justice Black)}

Justice Black concurred in the result and joined Parts First and Second of Stone's opinion, but expressly withheld his vote from Part Third (\href{https://supreme.justia.com/cases/federal/us/304/144/}{304 U.S. at 155}). Although Black did not write a separate opinion, his refusal to join Part Third reflected his skepticism of judicial fact-finding into the rationality of economic legislation (which Part Third still left open for as-applied challenges) as well as his textualist discomfort with open-ended judicial scrutiny untethered from explicit constitutional text.

\subsection{Concurrence in the Result (Justice Butler)}

Justice Butler concurred only in the result (\href{https://supreme.justia.com/cases/federal/us/304/144/}{304 U.S. at 155}). He agreed that the indictment stated a prima facie violation sufficient to defeat a demurrer, but insisted that the company must be permitted at trial to present evidence demonstrating that Congress's declarations of health danger and fraud had no substantial foundation. Invoking the canon of constitutional avoidance, Butler warned that if the statute were construed to bar ``wholesome food products that demonstrably are neither injurious to health nor calculated to deceive,'' it would violate the Fifth Amendment (\href{https://supreme.justia.com/cases/federal/us/304/144/}{304 U.S. at 155}).

\subsection{Dissent (Justice McReynolds)}

Justice McReynolds dissented without written opinion, stating simply that the judgment sustaining the demurrer should be affirmed (\href{https://supreme.justia.com/cases/federal/us/304/144/}{304 U.S. at 155}). Consistent with his \textit{Lochner}-era jurisprudence, McReynolds viewed the Filled Milk Act as an unconstitutional federal intrusion on economic liberty and state police power.

\section{Examples: Future Applications}

\begin{enumerate}
    \item \textbf{Same-Side (Upheld under Rational Basis):} \textit{Occupational Licensing of Interior Designers.} A state requires commercial interior designers to complete a four-year degree and pass a state licensing exam, ostensibly to protect building safety and consumers, even though critics argue it primarily shields incumbent firms from competition. Under \textit{Carolene Products}, this is ordinary commercial regulation implicating neither a specific Bill of Rights prohibition nor a discrete and insular minority. Because the safety rationale is ``at least debatable,'' a court will presume supporting facts exist and uphold the law.
    \item \textbf{Same-Side (Upheld under Rational Basis):} \textit{Selective Ban on Plant-Based ``Butter'' Substitutes.} Congress prohibits interstate shipment of a newly formulated canola-oil spread packaged to resemble dairy butter, while leaving traditional margarine unregulated. Just as Congress could target filled milk without banning oleomargarine, a legislature may address commercial problems piecemeal; courts will not second-guess whether labeling rules would have sufficed.
    \item \textbf{Opposite-Side (Heightened Scrutiny / Invalidated):} \textit{State Ban on Distributing Political Pamphlets Without a Mayor's Permit.} A municipality prohibits handing out leaflets criticizing local zoning laws without prior official approval. Under Footnote Four's first and second paragraphs, the ordinance both falls within a specific First Amendment prohibition (incorporated via the Fourteenth Amendment) and clogs ``those political processes which can ordinarily be expected to bring about repeal of undesirable legislation.'' The presumption of constitutionality narrows, triggering strict scrutiny.
    \item \textbf{Opposite-Side (Heightened Scrutiny / Invalidated):} \textit{Discriminatory Exclusion of a Religious Minority from Public Contracting.} A state bars businesses owned by members of a small, unpopular religious sect from bidding on public construction contracts, citing generalized ``community trust'' concerns. Under Footnote Four's first and third paragraphs, this targets a discrete and insular religious minority unable to rely on ordinary majoritarian coalition-building, warranting searching judicial review and invalidation.
    \item \textbf{Fence-Sitter (Borderline / Contested):} \textit{Total Prohibition on a Safe, Vitamin-Fortified Food Substitute Backed Purely by Incumbent Lobbying.} Suppose a state completely bans the sale of lab-grown dairy protein that is nutritionally identical to cow's milk, clearly labeled, and certified safe by federal regulators, after legislative hearings reveal only testimony from local dairy farmers seeking to eliminate competition. Under Stone's deferential rational-basis test, the state would argue that economic protection of local agriculture or speculative long-term health caution supplies a conceivable rational basis. Yet under Part Third's caveat for as-applied proof of changed facts and Butler's concurrence---as well as modern ``rational basis with bite'' arguments---a challenger would contend that known scientific facts preclude any non-arbitrary public-interest basis.
\end{enumerate}

\section{Critique}

\textbf{Originalist, Textualist, and Libertarian Critiques:} Originalist and libertarian scholars criticize \textit{Carolene Products} from two distinct angles. First, scholars such as Randy Barnett and Evan Bernick argue that the Court's extreme ``presumption of constitutionality'' for economic legislation abdicates the judiciary's original duty under the Due Process Clauses to guard against arbitrary exercises of legislative power that serve private rent-seeking rather than constitutionally proper public ends (\href{https://scholarship.law.georgetown.edu/cgi/viewcontent.cgi?article=3162&context=facpub}{Barnett \& Bernick, \textit{No Arbitrary Power}}). Public-choice scholarship famously reinforced this critique: Geoffrey P.~Miller's landmark study, \textit{The True Story of Carolene Products}, 1987 Sup.~Ct.~Rev.~397, documented that the Filled Milk Act was largely special-interest legislation engineered by the powerful dairy lobby to crush a wholesome, lower-cost competitor (\href{https://chicagounbound.uchicago.edu/cgi/viewcontent.cgi?article=2988&context=journal_articles}{Strauss, \textit{Is Carolene Products Obsolete?}}; \href{https://openyls.law.yale.edu/bitstreams/3c63a79b-a00d-4a09-8de5-9363a63ca8a3/download}{Yale Law Scholarship}). Second, conservative textualists question the two-tiered hierarchy created by Footnote Four, arguing that the constitutional text nowhere elevates certain unenumerated personal liberties above economic and property rights, and that Footnote Four invited courts to construct an extratextual taxonomy of preferred freedoms (\href{https://ecollections.law.fiu.edu/cgi/viewcontent.cgi?article=1300&context=lawreview}{Colby \& Smith, discussed in FIU L.~Rev.}).

\textbf{Progressive and Doctrinal Perspectives:} Conversely, many progressive and process-oriented scholars celebrated Footnote Four as an ingenious resolution to the ``countermajoritarian difficulty''---the tension between judicial review and democratic self-government. Most notably, John Hart Ely's \textit{Democracy and Distrust} (1980) built his ``representation-reinforcement'' theory directly on Footnote Four's second and third paragraphs, arguing that unelected judges act most legitimately when clearing blockages in democratic channels and protecting groups excluded from normal political coalitions (\href{https://legaltheorylexicon.com/2026/06/legal-theory-lexicon-constitutional-theory/}{Legal Theory Lexicon}; \href{https://chicagounbound.uchicago.edu/cgi/viewcontent.cgi?article=2224&context=public_law_and_legal_theory}{Stephanopoulos, \textit{The Anti-Carolene Court}}). Yet even sympathetic scholars have identified major flaws in Footnote Four's sociology of power: Bruce Ackerman (\textit{Beyond Carolene Products}, 98 Harv.~L.~Rev.~713 (1985)) observed that ``discrete and insular'' minorities are sometimes \textit{better} positioned to organize politically than diffuse, anonymous groups (such as consumers or the poor), while David Strauss and Nicholas Stephanopoulos note that modern Supreme Court doctrine has drifted away from Footnote Four by applying strict scrutiny symmetrically to majorities and minorities alike (\href{https://chicagounbound.uchicago.edu/cgi/viewcontent.cgi?article=2988&context=journal_articles}{Strauss}; \href{https://chicagounbound.uchicago.edu/cgi/viewcontent.cgi?article=2224&context=public_law_and_legal_theory}{Stephanopoulos}; \href{https://www.californialawreview.org/print/rethinking-political-power-in-judicial-review}{Calif.~L.~Rev.}).

\textbf{Underlying Competing Values:} \textit{Carolene Products} sits at the crossroads of \textbf{democratic majoritarianism} (leaving economic and social policy to elected legislatures rather than repeating \textit{Lochner}-era judicial vetoing) versus \textbf{constitutional limits on faction and minority oppression} (ensuring that legislative majorities cannot entrench themselves, silence dissent, or target vulnerable groups).

\section{Key Quotations}

\begin{quote}
``Even in the absence of such aids, the existence of facts supporting the legislative judgment is to be presumed, for regulatory legislation affecting ordinary commercial transactions is not to be pronounced unconstitutional unless, in the light of the facts made known or generally assumed, it is of such a character as to preclude the assumption that it rests upon some rational basis within the knowledge and experience of the legislators.''
\end{quote}
\noindent --- Justice Harlan Fiske Stone, Opinion of the Court (\href{https://supreme.justia.com/cases/federal/us/304/144/}{304 U.S. at 152}; \href{https://www.courtlistener.com/opinion/103017/united-states-v-carolene-products-co/}{CourtListener}).

\begin{quote}
``The Fifth Amendment has no equal protection clause, and even that of the Fourteenth, applicable only to the states, does not compel their legislatures to prohibit all like evils, or none. A legislature may hit at an abuse which it has found, even though it has failed to strike at another.''
\end{quote}
\noindent --- Justice Harlan Fiske Stone, Opinion of the Court (\href{https://supreme.justia.com/cases/federal/us/304/144/}{304 U.S. at 151}).

\begin{quote}
``There may be narrower scope for operation of the presumption of constitutionality when legislation appears on its face to be within a specific prohibition of the Constitution, such as those of the first ten amendments, which are deemed equally specific when held to be embraced within the Fourteenth.''
\end{quote}
\noindent --- Justice Harlan Fiske Stone, Footnote Four, Paragraph~1 (\href{https://supreme.justia.com/cases/federal/us/304/144/}{304 U.S. at 152 n.4}).

\begin{quote}
``Nor need we enquire whether similar considerations enter into the review of statutes directed at particular religious, \dots{} or national, \dots{} or racial minorities, \dots{}: whether prejudice against discrete and insular minorities may be a special condition, which tends seriously to curtail the operation of those political processes ordinarily to be relied upon to protect minorities, and which may call for a correspondingly more searching judicial inquiry.''
\end{quote}
\noindent --- Justice Harlan Fiske Stone, Footnote Four, Paragraph~3 (\href{https://supreme.justia.com/cases/federal/us/304/144/}{304 U.S. at 152--53 n.4}).

\begin{quote}
``If construed to exclude from interstate commerce wholesome food products that demonstrably are neither injurious to health nor calculated to deceive, they are repugnant to the Fifth Amendment.''
\end{quote}
\noindent --- Justice Pierce Butler, Concurring in the Result (\href{https://supreme.justia.com/cases/federal/us/304/144/}{304 U.S. at 155}).

\end{document}
